% AMC 12A 2023 Problem 15 Strategic Analysis
% Generated using Strategic Problem Solving Framework v2.1

\documentclass[11pt]{article}
\usepackage{amsmath, amssymb, amsthm}
\usepackage{geometry}
\usepackage{enumitem}
\usepackage{tcolorbox}
\tcbuselibrary{skins,breakable}
\usepackage{xcolor}

% Page setup
\geometry{margin=1in}
\setlength{\parindent}{0pt}
\setlength{\parskip}{6pt}

% Custom colored boxes
\newtcolorbox{problembox}[1][]{
    colback=blue!5!white,
    colframe=blue!75!black,
    title=Problem Configuration Analysis,
    breakable,
    #1
}

\newtcolorbox{strategybox}[1][]{
    colback=pink!5!white,
    colframe=pink!75!black,
    title=Strategic Framework Application,
    breakable,
    #1
}

\newtcolorbox{tacticalbox}[1][]{
    colback=green!5!white,
    colframe=green!75!black,
    title=Tactical Selection,
    breakable,
    #1
}

\newtcolorbox{conversionbox}[1][]{
    colback=yellow!5!white,
    colframe=yellow!75!black,
    title=Conversion Techniques,
    breakable,
    #1
}

\newtcolorbox{verificationbox}[1][]{
    colback=orange!5!white,
    colframe=orange!75!black,
    title=Verification Strategy,
    breakable,
    #1
}

\newtcolorbox{roadmapbox}[1][]{
    colback=purple!5!white,
    colframe=purple!75!black,
    title=Solution Roadmap,
    breakable,
    #1
}

\newtcolorbox{competitionbox}[1][]{
    colback=red!5!white,
    colframe=red!75!black,
    title=Competition Strategy,
    breakable,
    #1
}

\newtcolorbox{pitfallbox}[1][]{
    colback=gray!5!white,
    colframe=gray!75!black,
    title=Common Pitfalls and Warnings,
    breakable,
    #1
}

\newtcolorbox{solutionbox}[1][]{
    colback=cyan!5!white,
    colframe=cyan!75!black,
    title=Detailed Solution,
    breakable,
    #1
}

\newtcolorbox{keyinsightbox}[1][]{
    colback=magenta!5!white,
    colframe=magenta!75!black,
    title=Key Insights,
    breakable,
    #1
}

\begin{document}

\title{\textbf{AMC 12A 2023 Problem 15}\\Strategic Analysis}
\author{Using Framework v2.1}
\date{}
\maketitle

\section*{Problem Statement}

Usain is walking for exercise by zigzagging across a $100$-meter by $30$-meter rectangular field, beginning at point $A$ and ending on the segment $\overline{BC}$. He wants to increase the distance walked by zigzagging as shown in the figure. What angle $\theta = \angle PAB = \angle QPC = \angle RQB = \cdots$ will produce a length that is $120$ meters?

\[\textbf{(A) } \arccos\frac{5}{6} \qquad \textbf{(B) } \arccos\frac{4}{5} \qquad \textbf{(C) } \arccos\frac{3}{10} \qquad \textbf{(D) } \arcsin\frac{4}{5} \qquad \textbf{(E) } \arcsin\frac{5}{6}\]

\textit{Note: The figure shows a rectangular field $ABCD$ with $A$ at bottom-left, $B$ at bottom-right, $C$ at top-right, $D$ at top-left. The zigzag path goes from $A$ to $P$ (on $\overline{DC}$), then to $Q$ (on $\overline{AB}$), then to $R$ (on $\overline{DC}$), and so on, with all angles $\theta$ equal.}

\begin{problembox}
\section*{1. Problem Configuration Analysis}

\subsection*{A. Problem Classification}
\begin{itemize}[leftmargin=*]
    \item \textbf{Competition}: AMC 12A 2023
    \item \textbf{Problem Number}: 15
    \item \textbf{Difficulty Band}: Mid-range (Problem 15/25)
    \item \textbf{Expected Time}: 3-4 minutes
    \item \textbf{Point Value}: 6 points (standard AMC 12 scoring)
\end{itemize}

\subsection*{B. Mathematical Domain Classification}
\begin{itemize}[leftmargin=*]
    \item \textbf{Primary Domain}: Geometry / Trigonometry
    \item \textbf{Secondary Domain}: Path Optimization / Reflection
    \item \textbf{Tertiary Domain}: Similar Triangles / Pythagorean Theorem
    \item \textbf{Cross-Domain Potential}: High (connects geometry, trigonometry, physics concepts)
\end{itemize}

\subsection*{C. Problem Type}
\begin{itemize}[leftmargin=*]
    \item \textbf{Format}: Multiple choice (angle determination)
    \item \textbf{Question Type}: ``What angle...'' (inverse trigonometric function)
    \item \textbf{Core Task}: Find angle $\theta$ that produces a specific path length
    \item \textbf{Answer Form}: Inverse trigonometric expression
\end{itemize}

\subsection*{D. Given Information}
\begin{itemize}[leftmargin=*]
    \item Rectangular field dimensions: $100$ m (length) $\times$ $30$ m (width)
    \item Starting point: $A$ (corner of rectangle)
    \item Ending location: somewhere on segment $\overline{BC}$
    \item All zigzag angles are equal: $\theta = \angle PAB = \angle QPC = \angle RQB = \cdots$
    \item Total path length: $120$ meters
    \item The number of zigzag segments is not specified
\end{itemize}

\subsection*{E. Constraints}
\begin{itemize}[leftmargin=*]
    \item Path must zigzag between top and bottom edges
    \item All angles $\theta$ must be equal
    \item Path starts at $A$ and ends on $\overline{BC}$
    \item Total length must be exactly $120$ meters
    \item Field dimensions are fixed
\end{itemize}

\subsection*{F. Objective}
\begin{itemize}[leftmargin=*]
    \item \textbf{Find}: The angle $\theta$ that produces a $120$-meter path
    \item \textbf{Output}: Select the correct inverse trigonometric expression
\end{itemize}

\subsection*{G. Key Structural Features}
\begin{itemize}[leftmargin=*]
    \item \textbf{Reflection property}: Zigzag path resembles light reflecting between mirrors
    \item \textbf{Unfolding potential}: Path can be ``unfolded'' into a straight line
    \item \textbf{Similar triangles}: All zigzag segments form similar right triangles
    \item \textbf{Constant ratio}: Horizontal distance : path length = $100:120 = 5:6$
    \item \textbf{Physics analogy}: Horizontal velocity component is constant
\end{itemize}

\subsection*{H. Strategic Orientation}
\begin{itemize}[leftmargin=*]
    \item \textbf{Primary Strategy}: Unfolding technique (flatten zigzag into straight line)
    \item \textbf{Key Insight}: $\cos\theta$ represents horizontal progress per unit path length
    \item \textbf{Expected Difficulty}: Moderate - requires geometric insight or trigonometric setup
    \item \textbf{Time Pressure}: Should be completed in 3-4 minutes
\end{itemize}
\end{problembox}

\begin{strategybox}
\section*{2. Strategic Framework Application}

\subsection*{A. Psychological Strategies}
\begin{itemize}[leftmargin=*]
    \item \textbf{Confidence Calibration}: 
    \begin{itemize}
        \item Mid-range problem with elegant solution
        \item Multiple solution approaches available
        \item Answer choices help narrow possibilities
    \end{itemize}
    \item \textbf{Pattern Recognition}:
    \begin{itemize}
        \item Recognize zigzag/reflection pattern
        \item Connection to unfolding technique
        \item Ratio $100:120 = 5:6$ suggests $\cos\theta = 5/6$
    \end{itemize}
    \item \textbf{Problem Familiarity}:
    \begin{itemize}
        \item Similar to light reflection problems
        \item Related to path optimization
        \item Common geometry/trig problem type
    \end{itemize}
\end{itemize}

\subsection*{B. Investigation Strategies}
\begin{itemize}[leftmargin=*]
    \item \textbf{Initial Exploration}:
    \begin{itemize}
        \item Observe that horizontal distance is $100$ m, path length is $120$ m
        \item Ratio: $\frac{100}{120} = \frac{5}{6}$
        \item This suggests $\cos\theta = 5/6$ (horizontal component)
    \end{itemize}
    \item \textbf{Visualization Techniques}:
    \begin{itemize}
        \item Draw one zigzag segment with its projection
        \item Imagine unfolding the path into a straight line
        \item Consider the physics analogy: constant horizontal velocity
    \end{itemize}
    \item \textbf{Multiple Approaches}:
    \begin{itemize}
        \item Method 1: Unfolding (most elegant)
        \item Method 2: Similar triangles and proportions
        \item Method 3: Algebraic setup with Pythagorean theorem
        \item Method 4: Physics/velocity decomposition
    \end{itemize}
\end{itemize}

\subsection*{C. Argument Construction}
\begin{itemize}[leftmargin=*]
    \item \textbf{Logical Structure (Unfolding Method)}:
    \begin{enumerate}
        \item Visualize unfolding the zigzag path
        \item When unfolded, path becomes a straight line of length $120$ m
        \item Horizontal distance remains $100$ m
        \item In the right triangle formed: $\cos\theta = \frac{100}{120} = \frac{5}{6}$
        \item Therefore: $\theta = \arccos\frac{5}{6}$
    \end{enumerate}
    \item \textbf{Key Deductions}:
    \begin{itemize}
        \item Horizontal progress per unit path length = $\cos\theta$
        \item Total horizontal progress = (total path length) $\times$ $\cos\theta$
        \item Therefore: $120 \cos\theta = 100$
        \item So: $\cos\theta = 5/6$
    \end{itemize}
\end{itemize}

\subsection*{D. Cross-Domain Translation}
\begin{itemize}[leftmargin=*]
    \item \textbf{Geometric Interpretation}: Unfolding creates a right triangle
    \item \textbf{Trigonometric Interpretation}: $\cos\theta$ is the horizontal component
    \item \textbf{Physics Interpretation}: Horizontal velocity component is constant
    \item \textbf{Algebraic Interpretation}: Solve for $\theta$ using constraint equation
\end{itemize}
\end{strategybox}

\begin{tacticalbox}
\section*{3. Tactical Methodology Selection}

\subsection*{A. Core Mathematical Tactics}
\begin{itemize}[leftmargin=*]
    \item \textbf{Unfolding/Reflection}: Flatten zigzag path into straight line
    \item \textbf{Similar Triangles}: All zigzag segments are similar
    \item \textbf{Trigonometric Ratios}: Use $\cos\theta$ for horizontal component
    \item \textbf{Ratio Analysis}: Compare horizontal distance to path length
    \item \textbf{Pythagorean Theorem}: Relate segment lengths (alternative approach)
\end{itemize}

\subsection*{B. Domain-Specific Tactics}
\begin{itemize}[leftmargin=*]
    \item \textbf{Geometry}:
    \begin{itemize}
        \item Unfolding preserves path length but changes configuration
        \item Reflection angles create equal angles
        \item Similar triangles have proportional sides
    \end{itemize}
    \item \textbf{Trigonometry}:
    \begin{itemize}
        \item $\cos\theta$ represents horizontal progress rate
        \item Inverse trig functions recover angle from ratio
        \item Answer choices indicate $\arccos$ vs $\arcsin$
    \end{itemize}
    \item \textbf{Problem-Solving Techniques}:
    \begin{itemize}
        \item Draw auxiliary lines (projections)
        \item Use physics intuition (velocity components)
        \item Check answer choices for plausibility
    \end{itemize}
\end{itemize}

\subsection*{C. Tactical Priorities}
\begin{enumerate}[leftmargin=*]
    \item Recognize the $100:120 = 5:6$ ratio (30 seconds)
    \item Connect ratio to $\cos\theta$ (45 seconds)
    \item Visualize unfolding or similar triangles (60 seconds)
    \item Set up equation: $\cos\theta = 5/6$ (30 seconds)
    \item Identify answer: $\arccos(5/6)$ (30 seconds)
    \item Verify with answer choices (30 seconds)
\end{enumerate}
\end{tacticalbox}

\begin{conversionbox}
\section*{4. Conversion Techniques Application}

\subsection*{A. Method 1: Unfolding (Most Elegant)}
\begin{itemize}[leftmargin=*]
    \item \textbf{Concept}: Imagine ``unfolding'' the zigzag by reflecting alternating segments
    \item When unfolded, the path becomes a straight line from $A$ to some point $S'$
    \item The unfolded path has:
    \begin{itemize}
        \item Horizontal distance: $100$ m (same as rectangle length)
        \item Total path length: $120$ m (given)
    \end{itemize}
    \item This creates a right triangle where:
    \begin{align*}
    \cos\theta &= \frac{\text{adjacent (horizontal)}}{\text{hypotenuse (path)}}\\
    &= \frac{100}{120}\\
    &= \frac{5}{6}
    \end{align*}
    \item Therefore: $\theta = \arccos\frac{5}{6}$
\end{itemize}

\subsection*{B. Method 2: Similar Triangles and Proportions}
\begin{itemize}[leftmargin=*]
    \item Drop perpendicular from $P$ to $\overline{AB}$, meeting at point $Q'$
    \item Let $AQ' = x$, then $Q'P = 30$ (width of rectangle)
    \item By Pythagorean theorem: $AP = \sqrt{x^2 + 900}$
    \item All zigzag segments are similar, so they're all proportional
    \item Total horizontal distance = sum of all horizontal projections = $100$
    \item Total path length = sum of all segment lengths = $120$
    \item The ratio is preserved: if horizontal is $x$, path is $\frac{6x}{5}$
    \item Therefore: $\cos\theta = \frac{x}{\frac{6x}{5}} = \frac{5}{6}$
\end{itemize}

\subsection*{C. Method 3: Algebraic Setup}
\begin{itemize}[leftmargin=*]
    \item Let $x = $ horizontal projection of one full segment
    \item Then the hypotenuse of one full segment = $\sqrt{x^2 + 900}$
    \item Number of segments = $\frac{100}{x}$ (may not be integer, but proportions work)
    \item Total path length:
    \[\frac{100}{x} \cdot \sqrt{x^2 + 900} = 120\]
    \item Simplify:
    \[\sqrt{x^2 + 900} = \frac{6x}{5}\]
    \item Square both sides:
    \[x^2 + 900 = \frac{36x^2}{25}\]
    \[25x^2 + 22500 = 36x^2\]
    \[11x^2 = 22500\]
    \[x^2 = \frac{22500}{11}\]
    \item Then:
    \[\cos\theta = \frac{x}{\sqrt{x^2 + 900}} = \frac{x}{\frac{6x}{5}} = \frac{5}{6}\]
\end{itemize}

\subsection*{D. Method 4: Physics Analogy}
\begin{itemize}[leftmargin=*]
    \item Imagine Usain walks at constant speed
    \item Horizontal velocity component is constant (like light reflecting off mirrors)
    \item Horizontal velocity = (total velocity) $\times$ $\cos\theta$
    \item Time is proportional to distance walked
    \item Horizontal distance = (horizontal velocity) $\times$ (time)
    \item So: horizontal distance = (total distance) $\times$ $\cos\theta$
    \item Therefore: $100 = 120 \cos\theta$
    \item Thus: $\cos\theta = \frac{100}{120} = \frac{5}{6}$
\end{itemize}

\subsection*{E. Answer Selection}
\begin{itemize}[leftmargin=*]
    \item We found $\cos\theta = \frac{5}{6}$
    \item Therefore $\theta = \arccos\frac{5}{6}$
    \item Check answer choices: This is \textbf{(A)}
    \item Verify it's not $\arcsin$: $\sin\theta = \sqrt{1 - \cos^2\theta} = \sqrt{1 - \frac{25}{36}} = \sqrt{\frac{11}{36}} = \frac{\sqrt{11}}{6} \neq \frac{5}{6}$
\end{itemize}
\end{conversionbox}

\begin{verificationbox}
\section*{5. Verification Strategy Design}

\subsection*{A. Verification Level}
\begin{itemize}[leftmargin=*]
    \item \textbf{Selected Level}: Level 3 - Reasonableness Check + Alternative Method
    \item \textbf{Reasoning}: Geometric problem with elegant solution
    \item \textbf{Time Budget}: 45-60 seconds
\end{itemize}

\subsection*{B. Verification Checklist}
\begin{itemize}[leftmargin=*]
    \item \textbf{Check ratio}:
    \begin{itemize}
        \item $\frac{100}{120} = \frac{5}{6}$ $\checkmark$
        \item This matches $\cos\theta = 5/6$
    \end{itemize}
    \item \textbf{Verify angle type}:
    \begin{itemize}
        \item Since $\cos\theta = 5/6 < 1$, angle exists $\checkmark$
        \item Since $5/6 > 0$, angle is acute $\checkmark$
        \item Answer should be $\arccos$, not $\arcsin$ $\checkmark$
    \end{itemize}
    \item \textbf{Check answer format}:
    \begin{itemize}
        \item Answer (A): $\arccos\frac{5}{6}$ matches our result $\checkmark$
    \end{itemize}
    \item \textbf{Reasonableness}:
    \begin{itemize}
        \item Path length $120 > 100$ (straight distance) makes sense $\checkmark$
        \item Zigzag adds $20\%$ extra distance, reasonable $\checkmark$
    \end{itemize}
\end{itemize}

\subsection*{C. Alternative Verification Methods}
\begin{itemize}[leftmargin=*]
    \item \textbf{Method 1}: Compute $\theta \approx \arccos(0.833) \approx 33.6^\circ$, seems reasonable for zigzag
    \item \textbf{Method 2}: Check that $\sin\theta = \frac{\sqrt{11}}{6}$ gives vertical component
    \item \textbf{Method 3}: Verify with one specific segment using Pythagorean theorem
\end{itemize}

\subsection*{D. Answer Choice Validation}
\begin{itemize}[leftmargin=*]
    \item Our answer: $\arccos\frac{5}{6}$
    \item Matches choice: \textbf{(A)}
    \item Why not the others?
    \begin{itemize}
        \item (B) $\arccos\frac{4}{5}$: Would give $100/120 = 4/5$, but that's $0.8 \neq 5/6 \approx 0.833$
        \item (C) $\arccos\frac{3}{10}$: Would need $\cos\theta = 0.3$, too small (very large angle)
        \item (D) $\arcsin\frac{4}{5}$: Uses $\arcsin$ instead of $\arccos$
        \item (E) $\arcsin\frac{5}{6}$: Uses $\arcsin$ instead of $\arccos$
    \end{itemize}
\end{itemize}
\end{verificationbox}

\begin{roadmapbox}
\section*{6. Solution Roadmap}

\subsection*{A. High-Level Solution Path}
\begin{enumerate}[leftmargin=*]
    \item \textbf{Recognize pattern}: Horizontal distance $100$ m, path length $120$ m
    \item \textbf{Identify key ratio}: $\frac{100}{120} = \frac{5}{6}$
    \item \textbf{Connect to trigonometry}: This ratio represents $\cos\theta$
    \item \textbf{Solve for angle}: $\theta = \arccos\frac{5}{6}$
    \item \textbf{Select answer}: Choice (A)
\end{enumerate}

\subsection*{B. Detailed Solution Steps}

\textbf{Step 1}: Understand the geometry.
\begin{itemize}
    \item Rectangle is $100$ m $\times$ $30$ m
    \item Zigzag path starts at $A$, ends somewhere on $\overline{BC}$
    \item Total path length is $120$ m
    \item All angles $\theta$ are equal
\end{itemize}

\textbf{Step 2}: Use the unfolding technique.
\begin{itemize}
    \item Imagine ``unfolding'' the zigzag path by reflecting segments
    \item When unfolded, the path becomes a straight line
    \item This straight line has length $120$ m
    \item The horizontal distance covered is $100$ m (rectangle length)
\end{itemize}

\textbf{Step 3}: Apply trigonometry.
\begin{itemize}
    \item The unfolded path forms the hypotenuse of a right triangle
    \item The horizontal distance is the adjacent side
    \item Therefore: $\cos\theta = \frac{\text{horizontal}}{\text{path}} = \frac{100}{120} = \frac{5}{6}$
\end{itemize}

\textbf{Step 4}: Solve for the angle.
\begin{itemize}
    \item $\theta = \arccos\frac{5}{6}$
\end{itemize}

\textbf{Step 5}: Select the answer.
\begin{itemize}
    \item This matches choice \textbf{(A)}
\end{itemize}

\subsection*{C. Key Decision Points}
\begin{itemize}[leftmargin=*]
    \item \textbf{Decision 1}: Recognize that horizontal distance / path length gives $\cos\theta$
    \item \textbf{Decision 2}: Use unfolding technique (or similar approach)
    \item \textbf{Decision 3}: Identify answer as $\arccos$, not $\arcsin$
\end{itemize}

\subsection*{D. Time Management}
\begin{itemize}[leftmargin=*]
    \item Understanding problem: 30 seconds
    \item Computing ratio $100/120 = 5/6$: 20 seconds
    \item Connecting to $\cos\theta$: 45 seconds
    \item Solving for $\theta$: 30 seconds
    \item Verification: 45 seconds
    \item \textbf{Total}: 3 minutes 30 seconds
\end{itemize}
\end{roadmapbox}

\begin{competitionbox}
\section*{7. Competition Strategy Integration}

\subsection*{A. Problem Triage}
\begin{itemize}[leftmargin=*]
    \item \textbf{Difficulty Assessment}: Medium (Problem 15/25)
    \item \textbf{Confidence Level}: High if you recognize the unfolding technique
    \item \textbf{Decision}: Attempt if comfortable with geometry/trig; otherwise return later
\end{itemize}

\subsection*{B. Time Allocation}
\begin{itemize}[leftmargin=*]
    \item \textbf{Target Time}: 3-4 minutes
    \item \textbf{Maximum Time}: 5 minutes
    \item \textbf{Quick Approach}: If you see $100:120 = 5:6$ immediately, guess \textbf{(A)} $\arccos\frac{5}{6}$
\end{itemize}

\subsection*{C. Solution Format}
\begin{itemize}[leftmargin=*]
    \item Multiple choice: Select \textbf{(A)} $\arccos\frac{5}{6}$
    \item Verification: Check ratio and answer format
    \item Bubble answer sheet carefully
\end{itemize}

\subsection*{D. Strategic Considerations}
\begin{itemize}[leftmargin=*]
    \item \textbf{Pattern Recognition}: If you've seen unfolding problems before, this is quick
    \item \textbf{Elimination}: Can eliminate (D) and (E) if you determine answer uses $\arccos$
    \item \textbf{Ratio Clue}: $5:6$ ratio strongly suggests answer (A)
    \item \textbf{Time Saver}: Don't overthink - the elegant solution is usually correct
\end{itemize}
\end{competitionbox}

\begin{pitfallbox}
\section*{Common Pitfalls and Warnings}

\subsection*{A. Conceptual Errors}
\begin{itemize}[leftmargin=*]
    \item \textbf{Confusing $\arcsin$ with $\arccos$}: The horizontal component uses cosine, not sine
    \item \textbf{Using wrong ratio}: Don't use $\frac{30}{120}$ or $\frac{30}{100}$ - these don't directly give the angle
    \item \textbf{Forgetting unfolding}: Trying to set up complex equations instead of using elegant unfolding
    \item \textbf{Assuming integer segments}: The number of zigzag segments doesn't need to be an integer
\end{itemize}

\subsection*{B. Computational Errors}
\begin{itemize}[leftmargin=*]
    \item \textbf{Ratio simplification}: $\frac{100}{120} = \frac{5}{6}$, not $\frac{4}{5}$ or other values
    \item \textbf{Angle type}: Answer is $\arccos$, not $\arcsin$ (check which trig function you're inverting)
    \item \textbf{Arithmetic}: $100 \div 120 = 5 \div 6$, verify this carefully
\end{itemize}

\subsection*{C. Strategic Errors}
\begin{itemize}[leftmargin=*]
    \item \textbf{Overcomplicating}: Don't set up unnecessary variables or equations
    \item \textbf{Ignoring answer choices}: The choices tell you the answer involves $\frac{5}{6}$, $\frac{4}{5}$, or $\frac{3}{10}$
    \item \textbf{Time management}: If stuck after 3 minutes, guess (A) and move on
\end{itemize}

\subsection*{D. Warning Signs}
\begin{itemize}[leftmargin=*]
    \item If you get $\sin\theta = \frac{5}{6}$, reconsider - horizontal component is cosine
    \item If you get $\cos\theta = \frac{4}{5}$, recheck your ratio calculation
    \item If your answer isn't among the choices, revisit the unfolding concept
\end{itemize}
\end{pitfallbox}

\begin{solutionbox}
\subsection*{A. Main Solution (Unfolding Method)}
\begin{enumerate}[leftmargin=*]
    \item \textbf{Visualize the unfolding}:
    \begin{itemize}
        \item The zigzag path reflects between top and bottom edges
        \item Imagine ``unfolding'' by reflecting alternating segments
        \item The unfolded path becomes a straight line from $A$ to endpoint
    \end{itemize}
    
    \item \textbf{Identify distances}:
    \begin{itemize}
        \item Horizontal distance (rectangle length): $100$ m
        \item Total path length (unfolded hypotenuse): $120$ m
    \end{itemize}
    
    \item \textbf{Apply trigonometry}:
    \begin{align*}
    \cos\theta &= \frac{\text{adjacent}}{\text{hypotenuse}}\\
    &= \frac{\text{horizontal distance}}{\text{path length}}\\
    &= \frac{100}{120}\\
    &= \frac{5}{6}
    \end{align*}
    
    \item \textbf{Solve for angle}:
    \[\theta = \arccos\frac{5}{6}\]
\end{enumerate}

\subsection*{B. Alternative Solution (Similar Triangles)}
\begin{enumerate}[leftmargin=*]
    \item Let $x$ be the horizontal projection of one zigzag segment
    \item The vertical component is $30$ m (width of rectangle)
    \item By Pythagorean theorem, segment length = $\sqrt{x^2 + 900}$
    \item All segments are similar (same angle $\theta$)
    \item Total horizontal distance: $100$ m
    \item Total path length: $120$ m
    \item Since all segments are similar with the same ratio:
    \[\frac{\text{total horizontal}}{\text{total path}} = \frac{x}{\sqrt{x^2 + 900}}\]
    \item This gives us:
    \[\frac{100}{120} = \cos\theta = \frac{5}{6}\]
    \item Therefore: $\theta = \arccos\frac{5}{6}$
\end{enumerate}

\subsection*{C. Physics Interpretation}
\begin{itemize}[leftmargin=*]
    \item Imagine Usain walks at constant speed
    \item The horizontal velocity component is constant (like light reflecting)
    \item Horizontal velocity = (total velocity) $\times$ $\cos\theta$
    \item Since time is the same, distance ratio equals velocity ratio:
    \[\frac{\text{horizontal distance}}{\text{total distance}} = \cos\theta = \frac{100}{120} = \frac{5}{6}\]
    \item Therefore: $\theta = \arccos\frac{5}{6}$
\end{itemize}

\subsection*{D. Verification}
\begin{itemize}[leftmargin=*]
    \item Check ratio: $\frac{100}{120} = \frac{5}{6}$ $\checkmark$
    \item Verify answer type: $\arccos$ (not $\arcsin$) $\checkmark$
    \item Matches choice (A): $\arccos\frac{5}{6}$ $\checkmark$
    \item Reasonableness: $\theta \approx 33.6^\circ$, appropriate for zigzag $\checkmark$
\end{itemize}
\end{solutionbox}

\begin{keyinsightbox}
\textbf{Key Insight}: The most elegant approach is the \textbf{unfolding technique}. When you ``unfold'' a zigzag path by reflecting alternating segments, the path becomes a straight line. The horizontal distance remains $100$ m, and the total path length is $120$ m. This creates a right triangle where:

\[\cos\theta = \frac{\text{horizontal distance}}{\text{path length}} = \frac{100}{120} = \frac{5}{6}\]

Therefore: $\theta = \arccos\frac{5}{6}$

\textbf{General Pattern}: For zigzag problems:
\begin{itemize}
    \item Horizontal progress = (path length) $\times$ $\cos\theta$
    \item Vertical progress = (path length) $\times$ $\sin\theta$
    \item Unfolding preserves lengths and reveals the geometry
\end{itemize}

\textbf{Physics Analogy}: Like light reflecting between parallel mirrors, the horizontal velocity component remains constant. This is why the ratio of horizontal distance to total path length equals $\cos\theta$.

\textbf{Why this works}: The unfolding technique works because:
\begin{itemize}
    \item Each reflection preserves the angle
    \item The sum of horizontal projections equals the rectangle length
    \item The sum of segment lengths equals the total path length
    \item These ratios are preserved in the unfolded configuration
\end{itemize}
\end{keyinsightbox}

\section*{Final Answer}

The angle $\theta$ that produces a path length of $120$ meters is:

\[\boxed{\textbf{(A) } \arccos\frac{5}{6}}\]

\textbf{Solution Summary}:
\begin{itemize}
    \item Horizontal distance: $100$ m
    \item Path length: $120$ m
    \item Using unfolding: $\cos\theta = \frac{100}{120} = \frac{5}{6}$
    \item Therefore: $\theta = \arccos\frac{5}{6}$
\end{itemize}

\section*{Confidence Assessment}
\begin{itemize}[leftmargin=*]
    \item \textbf{Confidence Level}: 99\%
    \begin{itemize}
        \item Elegant unfolding technique gives clear answer
        \item Simple ratio calculation: $100/120 = 5/6$
        \item Multiple verification methods confirm result
        \item Answer format matches problem type
    \end{itemize}
    \item \textbf{Verification Applied}: Level 3 - Reasonableness + Alternative Method
    \begin{itemize}
        \item Checked ratio calculation
        \item Verified answer type ($\arccos$ vs $\arcsin$)
        \item Confirmed with physics interpretation
        \item All methods agree
    \end{itemize}
    \item \textbf{Time-Confidence Trade-off}: High confidence achieved in 3-4 minutes
\end{itemize}

\end{document}

